\chapter{Experimental Design and Evaluation}
\label{chap:evaluation}

\section{Evaluation objectives}

Evaluation addresses six distinct questions: (1) whether the graph obeys its
structural contract; (2) whether source objects and graph evidence are bound to
the recorded inputs; (3) whether sampled accepted edges are semantically
supported; (4) how well the topic classifier agrees with the available reference
labels; (5) how much coverage is lost through unresolved identities and
affiliations; and (6) whether important network conclusions persist under
specified identity or support policies. These objectives use different
denominators and are not combined into one accuracy number.

\section{Structural integrity and source consistency}

The graph validator checks unique IDs, valid typed endpoints, allowed labels and
predicates, accepted status, evidence existence, and recomputed support/weight
rules for derived edges. The release reports all ten predicates as nonempty and
passes the stated size gates. Neo4j inventory, integrity queries, and replay
checks independently confirm that the stored graph matches the frozen artifact.
These checks establish internal consistency of the data representation, not
whether the source record or person identity is semantically correct.

The v4 source-population check compares retained author and reference objects
for all 13,606 matched works, covering 65,469 positions and 341 raw response
hashes, against the immutable OpenAlex cache. This is byte/source consistency.
Because the same provider records informed graph construction, it is not an
independent external validation of authorship, affiliation, or identity.

\section{Accepted-edge precision audits}

\subsection{Prospective v4 audit}

The prospective audit sampled accepted edges from the frozen v4 frame of
225,731 edges. It drew 200 uniform global edges and ten supplementary edges per
relation; the supplements were kept separate from the global estimate. The
result comprised 300 unique edges and 350 assignments, including 50 global
overlap edges judged by both evaluators. Two separate fresh GPT-6 review contexts
completed 175 assignments each. Their model/provider build identifiers were not
available. Separate contexts and nonshared answers do not establish statistical
independence.

The uniform sample received 193 correct/supportable, one incorrect, six
unverifiable, and zero missing judgments. Unverifiable is counted as incorrect
for the precision denominator, yielding $193/200=96.50\%$. The Wilson 95\%
interval is 92.95\%--98.29\%. There was one categorical disagreement among the
50 overlap edges before adjudication (49/50 agreement; Cohen's $\kappa=0.6599$).
The scoring artifact passes the project's AI-assessed point and interval
comparisons, but it does not certify the course's independently verified
accuracy requirement.

\begin{table}[tbp]
  \centering
  \caption{Global accepted-edge audit results, kept separate by graph version.}
  \label{tab:edge-audits}
  \small
  \begin{tabularx}{0.98\textwidth}{>{\raggedright\arraybackslash}p{2.2cm}>{\raggedright\arraybackslash}p{6cm}>{\centering\arraybackslash}p{3cm}>{\raggedright\arraybackslash}X}
    \toprule
    \textbf{Audit} & \textbf{Results} & \textbf{Wilson interval} & \textbf{Interpretation}\\
    \midrule
    Corrected v4 & 193/200 supported; 1 incorrect; 6 unverifiable & 92.95--98.29\% & AI-assessed precision; not human-verified\\
    Historical v3 & 176/200 supported; 23 incorrect; 1 unverifiable & 82.77--91.80\% & Historical result; numeric checks failed\\
    Sprint, ten relations & 76/80 supported; 3 unsupported; 1 unknown & Not pooled & Relation-stratified diagnostic only\\
    \bottomrule
  \end{tabularx}
  \begin{minipage}{0.98\textwidth}
    \vspace{0.35em}\footnotesize Unknown/unverifiable judgments count as nonsupport in precision. The first two rows sample global accepted-edge frames; the sprint row uses eight cases per relation and is not an estimate for the full graph.
  \end{minipage}
\end{table}

The ten-per-relation rows are supplementary diagnostics. Each has $n=10$ and
wide uncertainty; a perfect 10/10 row still has a Wilson lower bound of about
72.25\%. In these supplements, `CONTRIBUTES\_TO\_TOPIC` had 9/10 supported and
one incorrect; `RESEARCHES\_TOPIC` had 8/10 supported and two unverifiable;
the other eight relations were 10/10 under the AI judgments. These numbers are
not pooled with the uniform sample.

\begin{table}[!htbp]
  \centering
  \caption[V4 relation-level AI diagnostics]{V4 AI-assessed relation supplements. Each row is a separate $n=10$
  diagnostic, not part of the 200-edge global estimator.}
  \label{tab:v4-relation-supplement}
  \small
  \begin{tabularx}{\textwidth}{>{\ttfamily\raggedright\arraybackslash}p{3.8cm}rrrX}
    \toprule
    \textbf{Predicate} & \textbf{Supported} & \textbf{Incorrect} &
    \textbf{Unverifiable} & \textbf{Wilson 95\% interval}\\
    \midrule
    AFFILIATED\_WITH & 10/10 & 0 & 0 & 72.25--100.00\%\\
    AUTHORED & 10/10 & 0 & 0 & 72.25--100.00\%\\
    CITES & 10/10 & 0 & 0 & 72.25--100.00\%\\
    COAUTHORED\_WITH & 10/10 & 0 & 0 & 72.25--100.00\%\\
    COLLABORATED\_WITH & 10/10 & 0 & 0 & 72.25--100.00\%\\
    \makecell[l]{\texttt{CONTRIBUTES\_TO\_}\\\texttt{TOPIC}} & 9/10 & 1 & 0 & 59.58--98.21\%\\
    HAS\_TOPIC & 10/10 & 0 & 0 & 72.25--100.00\%\\
    PUBLISHED\_IN & 10/10 & 0 & 0 & 72.25--100.00\%\\
    RESEARCHES\_TOPIC & 8/10 & 0 & 2 & 49.02--94.33\%\\
    SUBTOPIC\_OF & 10/10 & 0 & 0 & 72.25--100.00\%\\
    \bottomrule
  \end{tabularx}
\end{table}


All 350 assignments reuse construction sources. Evidence designation includes
both primary-source and secondary-source-only packets, but a primary proceedings
byline does not authenticate the mapping from an OpenAlex ID to a real person.
Some aggregate support counts were checked only through bounded paper witnesses.
The audit assesses sampled accepted-edge precision and cannot estimate the
recall of missing edges. Its interval also excludes evaluator errors and
source/model dependence.

\subsection{Preserved historical v3 audit}

The prior frozen v3 graph, \texttt{release-9d4a6fb20f63}, remains intact as a
historical result. Its 200-edge uniform sample was assessed as 176/200 (88.00\%),
with 23 incorrect and one unverifiable judgment; its Wilson 95\% interval was
82.77\%--91.80\%. The course point comparison and project precision checks
failed. These numbers describe v3, not corrected production v4. The audit
identified a missing observed-ORCID conflict path, which led to a separately
versioned v4 correction. The new score cannot be read as an independently
verified causal improvement: the old review motivated the correction, the source
snapshot is reused, and the new sample was not an untouched independent study.

\subsection{Deadline-sprint ten-relation review}

A later bounded sprint assessed eight accepted edges in each of the ten
relations, for 80 assignments. Two fresh GPT-6-Luna contexts used high reasoning
effort and disjoint 40-case packets. They judged 76 cases supported, three
unsupported, and one unknown. The four flagged cases comprise two unsupported
`RESEARCHES\_TOPIC` edges, one unsupported `CONTRIBUTES\_TO\_TOPIC` edge, and
one unknown `AUTHORED` identity mapping. This is a small, deliberately
relation-stratified diagnostic with unequal population inclusion probabilities;
it has no pooled population-accuracy estimator. It also reuses construction
sources and remains AI-assessed, not human ground truth.

\begin{table}[tbp]
  \centering
  \caption{Deadline-sprint Luna diagnostic: eight sampled accepted edges per
  relation. Denominators and inclusion probabilities differ across relations.}
  \label{tab:sprint-relation-audit}
  \scriptsize
  \begin{adjustbox}{max width=\textwidth}
  \begin{tabular}{lrrrrrr}
    \toprule
    \textbf{Predicate} & \textbf{Frame} & \textbf{Inclusion prob.} &
    \textbf{Supported} & \textbf{Unsupported} & \textbf{Unknown} & \textbf{Wilson 95\% interval}\\
    \midrule
    AFFILIATED\_WITH & 38,664 & 0.00020691 & 8/8 & 0 & 0 & 67.6--100.0\%\\
    AUTHORED & 50,437 & 0.00015861 & 7/8 & 0 & 1 & 52.9--97.8\%\\
    CITES & 7,308 & 0.00109469 & 8/8 & 0 & 0 & 67.6--100.0\%\\
    COAUTHORED\_WITH & 80,482 & 0.00009940 & 8/8 & 0 & 0 & 67.6--100.0\%\\
    COLLABORATED\_WITH & 13,031 & 0.00061392 & 8/8 & 0 & 0 & 67.6--100.0\%\\
    CONTRIBUTES\_TO\_TOPIC & 2,332 & 0.00343053 & 7/8 & 1 & 0 & 52.9--97.8\%\\
    HAS\_TOPIC & 14,222 & 0.00056251 & 8/8 & 0 & 0 & 67.6--100.0\%\\
    PUBLISHED\_IN & 13,617 & 0.00058750 & 8/8 & 0 & 0 & 67.6--100.0\%\\
    RESEARCHES\_TOPIC & 5,622 & 0.00142298 & 6/8 & 2 & 0 & 40.9--92.9\%\\
    SUBTOPIC\_OF & 16 & 0.50000000 & 8/8 & 0 & 0 & 67.6--100.0\%\\
    \bottomrule
  \end{tabular}
  \end{adjustbox}
  \begin{minipage}{\textwidth}
    \vspace{0.35em}\footnotesize Wilson intervals are conditional and descriptive; no finite-population correction or multiplicity adjustment is applied. The relation strata are deliberately not pooled into an overall precision estimate.
  \end{minipage}
\end{table}


\section{Topic-classifier evaluation}

The topic evaluation uses 50 frozen test papers. Two historical AI contexts each
reviewed 30 papers with ten overlapping assignments; their labels are reference
judgments, not human gold. Each paper contributes total weight one, and paired
overlap labels split that paper's weight between reviewers. Precision, recall,
and F1 are computed over the explicit considered-label sets; coverage is reported
separately from label agreement. The test papers were already exposed in earlier
review work, and the v4 edge audit generated no new topic labels.

The embedding method has test coverage 45/50 (90\%) and micro precision, recall,
and F1 of 0.6863, 0.3535, and 0.4667. The phrase baseline has coverage 13/50
(26\%) and micro scores 0.9286, 0.1313, and 0.2301. The embedding method
increases coverage and recall, while the phrase baseline is more precise on its
smaller set of assigned labels. Neither reaches the 0.95 precision target.
Macro F1 is 0.3627 for embeddings and 0.1473 for the phrase baseline. These are
historical agreements with AI-assigned references, not independently verified
topic accuracy.

\begin{table}[tbp]
  \centering
  \caption{Historical topic agreement against AI-assigned reference labels.}
  \label{tab:topic-evaluation}
  \footnotesize
  \begin{tabularx}{\textwidth}{p{2.4cm}*{7}{>{\centering\arraybackslash}X}}
    \toprule
    \textbf{Classifier} & \makecell{Coverage} & \makecell{Micro\\P} & \makecell{Micro\\R} & \makecell{Micro\\F1} & \makecell{Macro\\P} & \makecell{Macro\\R} & \makecell{Macro\\F1}\\
    \midrule
    Phrase baseline & 13/50 (26\%) & 0.9286 & 0.1313 & 0.2301 & 0.4167 & 0.0943 & 0.1473\\
    Embedding & 45/50 (90\%) & 0.6863 & 0.3535 & 0.4667 & 0.4986 & 0.3066 & 0.3627\\
    \bottomrule
  \end{tabularx}
  \begin{minipage}{\textwidth}
    \vspace{0.35em}\footnotesize Coverage is the number of papers with at least one assigned topic. Test labels were already exposed in historical AI review; this is not an untouched human test set.
  \end{minipage}
\end{table}

\section{Identity and evidence sensitivity}

Sensitivity experiments compare network rankings under policy changes; they do
not label a candidate graph as more accurate. A strict evidence filter on v4
retains 7.78\% of researcher collaboration edges and 26.1\% of organization
collaboration edges. It removes citation and paper-topic edges because those
base predicates lack the aggregate support count used by that filter, while
retaining all R10 organization-topic edges because construction already requires
three supporting papers. These mechanical retention outcomes are not accuracy
estimates.

The separate Task2C study compares two experimental identity candidates with
frozen v4 on the same approximate-betweenness analysis. Degree top-20 membership
is stable (19/20 and 20/20 overlap), while betweenness top-20 overlap is 11/20
and 12/20. Candidate results remain development scenarios with unmeasured
semantic precision and recall. No v5 identity candidate is promoted.

\begin{table}[tbp]
  \centering
  \caption{Top-20 ranking overlap between frozen production v4 and experimental
  identity-policy candidates.}
  \label{tab:identity-sensitivity}
  \begin{tabular}{lrr}
    \toprule
    \textbf{Candidate versus v4} & \textbf{Degree overlap} & \textbf{Betweenness overlap}\\
    \midrule
    Candidate-v3 & 19/20 & 11/20\\
    Task2B & 20/20 & 12/20\\
    \bottomrule
  \end{tabular}
  \begin{minipage}{0.82\textwidth}
    \vspace{0.35em}\footnotesize These are sensitivity comparisons, not accuracy comparisons or graph promotions. Candidate identity precision and recall were not measured.
  \end{minipage}
\end{table}

\section{Metric definitions}

For a set of reviewed accepted edges, precision is the proportion of reviewed
claims judged supported:
\begin{equation}
  \widehat{P}=\frac{n_{\text{supported}}}{n_{\text{reviewed}}},
  \qquad
  n_{\text{reviewed}}=n_{\text{supported}}+n_{\text{unsupported}}+
  n_{\text{unknown}}+n_{\text{missing}}.
\end{equation}
The project counts unknown/unverifiable and missing items as nonsupport for the
precision calculation. Recall would require a defensible sample of absent as
well as present edges; the accepted-edge audit does not provide one. For topic
labels, $P=TP/(TP+FP)$, $R=TP/(TP+FN)$, and $F_1=2PR/(P+R)$, with micro scores
pooling decisions and macro scores averaging label-level results.

The Wilson interval for an observed proportion $\widehat{p}=x/n$ uses a normal
quantile $z$ and is
\begin{align}
  c &= \frac{\widehat{p}+z^2/(2n)}{1+z^2/n},\\
  h &= \frac{z}{1+z^2/n}
       \sqrt{\frac{\widehat{p}(1-\widehat{p})}{n}+\frac{z^2}{4n^2}},\\
  CI_{95\%} &= [c-h, c+h], \qquad z\approx1.96.
\end{align}
This interval describes sampling variability conditional on the judgments. It is
not an uncertainty interval for evaluator reliability or source correctness.

For a collaboration graph, degree is the number of distinct neighbors and
strength is the sum of incident fractional weights. Betweenness is normalized
in the analysis output and approximated with 256 pivots on weighted shortest
paths. Citation in-degree counts accepted incoming CITES edges from within the
corpus. None of these metrics is a universal measure of scientific influence.
